\scanpage{27}
$\gamma)$ \deleted{$T_0$} $X$ (pour sa relation d'ordre) est artinien, i.e. tout
\deleted{famille}\inserted{suite} décroissante d'éléments de $X$ est \unsure{station.}

$\delta)$ $\widetilde X$ est essentiel (i.e. chaque pt de $X$ est essentiel)

\margin{en fait, $\operatorname{Im}(X\to\widehat X)$ est l'ens des pts ess. de $\widetilde X$ !}

j) Soit $E$ un topos. Conditions équivalentes

$\alpha)$ $E$ \deleted{\unsure{isom.|n'est}}\inserted{équiv.} à un $\operatorname{Top}(X)$ \uline{et} à un $\widehat C$

$\beta)$ $E$ équiv. à $\operatorname{Top}(X)$, $X$ espace essentiel

$\gamma)$ $E$ équ. à $\widehat I$, $I$ ens. ordonné

$\delta)$ $E$ engendré par les sous-objets de l'objet final, et
a suff. de pts essentiels

[\deleted{\unsure{On}}\unsure{On} dit alors que $E$ ``est'' un espace top. essentiel.]

Prouver que si $C$ est une catégorie, $\widehat C$ satisfait aux conditions
précédentes sss $C$ est préordonnée.

Prouver que si $X$ est un espace top, $\operatorname{Top}(X)$ satisfait aux
conditions précédentes sss l'ens des $x\in X$ qui sont
essentiels est très dense.

\deleted{\unsure{k) Soit $X$}}
On a des équivalences de 2-catégories

\[
(\text{topos spatiaux essentiels})\mathrel{\overset{\approx}{\simeq}}
\deleted{\text{ens. ordonnés}}\inserted{(\mathrm{Préord})}
\mathrel{\overset{\approx}{\simeq}}(\mathrm{Esp.ess.})
\]

[NB. Si $X$ est \unsure{esp.|ess.} on le récupère à partir de $E=\operatorname{Top}(X)$ comme $\underline{\mathrm{Points}}(\unsure{E_{\mathrm p}|E})$]
% archive sheet number: 26
